% !PRINTED 450
\weberpage{450}
%<<CONT
\gesperrt{can be transformed into every other by substitutions of the group} $G_{168}$\gesperrt{.}

If we set
\weq{3}{T = x_1 + x_2 + x_3,}
then $T = 0$ is the equation of the line joining the points $\pi, \pi'$. The function $T$ remains unchanged by $\chi$ and changes its sign by $\omega$. But by the 28 substitutions $\Theta^s\tau^r$ it passes over into 28 \gesperrt{different} functions $T_{s,r}$, which we can easily represent, with the aid of the formulae of \S.\,115, (11), (14), as follows:
\weq{4}{
\begin{aligned}
T_{0,r} &= \varepsilon^r x_1 + \varepsilon^{2r} x_2 + \varepsilon^{4r} x_3 \\
h\,T_{1,r} &= \beta^2\varepsilon^{1+r} x_1 + \alpha^2\varepsilon^{2+2r} x_2 + \gamma^2\varepsilon^{4+4r} x_3 \\
-h\,T_{2,r} &= \alpha^2\varepsilon^{2+r} x_1 + \gamma^2\varepsilon^{4+2r} x_2 + \beta^2\varepsilon^{1+4r} x_3 \\
h\,T_{3,r} &= \gamma^2\varepsilon^{4+r} x_1 + \beta^2\varepsilon^{1+2r} x_2 + \alpha^2\varepsilon^{2+4r} x_3,
\end{aligned}
}
and these, set equal to zero, represent 28 different straight lines, which can all be derived from one of them by substitutions of the group.

On the line $T$ lie the three points with the coordinates
\[
\begin{aligned}
\alpha\gamma &: \beta\gamma : \beta\alpha \\
\beta\alpha &: \alpha\gamma : \beta\gamma \\
\beta\gamma &: \beta\alpha : \alpha\gamma,
\end{aligned}
\]
which arise from (10), \S.\,118 by the substitutions $\tau^{-4}, \tau^{-4}\chi, \tau^{-4}\chi^2$, and thus belong to the eight-fold poles.

6. \gesperrt{On every line $T$ lie three eight-fold and two three-fold poles.}

It still remains to investigate the third pole $\pi_0$ of the substitution $\chi$, with the coordinates
\[
x_1 = x_2 = x_3
\]
of which we have already proved that it remains unchanged by the six substitutions
\weq{5}{1, \chi, \chi^2; \omega, \omega\chi, \omega\chi^2}
The question now arises whether this pole is not more than six-fold.

If besides (5) there are still other substitutions which leave the pole $\pi_0$ unchanged, then these must necessarily be of degree $2$, since $\pi_0$ is neither a four-fold nor a seven-fold pole, and because, further, the degree of the group belonging to $\pi_0$ cannot be divisible by 9.
